\FloatBarrier
\section{Trabalhos relacionados}

M\'ascaras de inst\^ancia permitem avaliar fronteiras e separar objetos
adjacentes, mas exigem um ground truth ausente na CRIC. Pontos reduzem custo de
anota\c{c}\~ao e identificam inst\^ancias sem descrever sua extens\~ao
\cite{chen2021sparsepoints,xia2023colorization,chen2023p2bnet}. Em imagens
biom\'edicas, Chen et al. propuseram segmenta\c{c}\~ao histopatol\'ogica fracamente
supervisionada com pontos esparsos \cite{chen2021sparsepoints}, enquanto Xia et
al. exploraram pontos e coloriza\c{c}\~ao para segmenta\c{c}\~ao nuclear
\cite{xia2023colorization}. Na detec\c{c}\~ao de objetos, a rede Point-to-Box aprende a
extens\~ao a partir de anota\c{c}\~oes pontuais \cite{chen2023p2bnet}. Essa frente
fornece a ancoragem metodol\'ogica mais pr\'oxima deste artigo: o ponto \'e tratado
como supervis\~ao espacial incompleta, e a regi\~ao usada pelo detector passa a ser
uma escolha experimental.

Em citologia cervical, revis\~oes recentes destacam resultados elevados em
classifica\c{c}\~ao, mas tamb\'em grande heterogeneidade de bases, classes,
prepara\c{c}\~ao e protocolos de avalia\c{c}\~ao \cite{jiang2023systematic}. O
CerviCell-detector avaliou detectores de objetos em imagens de Papanicolau,
incluindo experimentos na CRIC, com formula\c{c}\~oes bin\'aria e Bethesda
\cite{kalbhor2023cervicell}. O presente estudo se aproxima dele pelo dom\'inio e
pela base de dados, mas remove a decis\~ao diagn\'ostica do alvo e pergunta se
todas as inst\^ancias pontuais podem ser localizadas e contadas. Assim, sua
contribui\c{c}\~ao compar\'avel \`a literatura n\~ao \'e superar numericamente
modelos de segmenta\c{c}\~ao ou classifica\c{c}\~ao, mas quantificar o efeito de
transformar um ponto em uma regi\~ao de treinamento para detectores de caixas.

Essa diferen\c{c}a tamb\'em orienta a compara\c{c}\~ao experimental. Trabalhos que
classificam c\'elulas ou les\~oes respondem a uma pergunta cl\'inica distinta,
enquanto m\'etodos de segmenta\c{c}\~ao dependem de contornos que n\~ao est\~ao
dispon\'iveis nas anota\c{c}\~oes usadas aqui. Por isso, este artigo evita uma
hierarquia num\'erica direta com resultados externos e privilegia a
reprodutibilidade do protocolo: origem dos dados, transforma\c{c}\~ao dos pontos,
crit\'erio de escolha da pseudo-caixa, particionamento por imagem e avalia\c{c}\~ao
cruzada final s\~ao descritos de forma expl\'icita.
